%%%%%%%%%%%%%%%%%%%%%%%%%%%%%%%%%%%%%%%%%%%%%%%%%%%%%%%%%%%%%%%%%%%%%%%%%%%%%%%
% Information measures (discrete)
%%%%%%%%%%%%%%%%%%%%%%%%%%%%%%%%%%%%%%%%%%%%%%%%%%%%%%%%%%%%%%%%%%%%%%%%%%%%%%%

\subsection{Surprisal and Entropy}
\label{sec:entropy-entropy}

\begin{definition}[Surprisal / self-information]
Let $X$ be a discrete random variable with pmf $p_X$ on a finite alphabet $\cX$.
The \emph{surprisal} (self-information) of an outcome $x\in\cX$ is
\begin{equation}
    S_X(x) \coloneqq \log\frac{1}{p_X(x)} = -\log p_X(x).
\end{equation}
\end{definition}

Rare events carry more information: if $p_X(x)$ is small, then $S_X(x)$ is large.

\begin{definition}[Entropy]
The (Shannon) entropy of $X$ is the expected surprisal:
\begin{equation}
    H(X)
    \coloneqq \E{S_X(X)}
    = \sum_{x\in\cX} p_X(x)\log\frac{1}{p_X(x)}
    = -\sum_{x\in\cX} p_X(x)\log p_X(x).
\end{equation}
\end{definition}

\paragraph{Interpretation.}
Entropy is the fundamental limit of lossless compression: roughly speaking, $H(X)$ is the minimum achievable average number of bits needed to describe one draw of $X$.
We make this precise in Section~\ref{sec:lossless-intro}.

\begin{example}[A highly biased coin]
Let $X\in\{0,1\}$ with $\Pr{X=1}=0.01$.
Then
\[
H(X)= -0.01\log 0.01 - 0.99\log 0.99 \approx 0.0808\ \text{bits}.
\]
Although there are two outcomes, the uncertainty is far below $1$ bit.
\end{example}

\subsubsection*{Basic properties}

\begin{theorem}[Jensen's Inequality]
Let $f$ be a concave function, then
\begin{align}
    \E{f(X)} \leq f(\E{X}).
\end{align}
\end{theorem}
Since the definition of concave function is that
\begin{align}
    \lambda f(x) + (1-\lambda) f(y) \leq f(\lambda x + (1-\lambda) y)
\end{align}
for all $x, y$ and $0\leq \lambda \leq 1$,
in other words, weighted average of $f$ is smaller or equal to $f$ at weighted averaged input.
Then, Jensen's inequality is a generalized version of it,
where $\E{f(X)}$ is weighted average of $f$ and $f(\E{X})$ is a function value at weighted average input.
For strictly concave function $f$, Jensen's inequality holds with equality if and only if $f(x)$ is constant.


\begin{theorem}[Non-negativity]
$H(X)\ge 0$, with equality if and only if $X$ is deterministic.
\end{theorem}
\begin{proof}
For $0<p\le 1$ we have $\log(1/p)\ge 0$.
Thus each term $p_X(x)\log(1/p_X(x))$ is nonnegative, so their sum is nonnegative.
If $H(X)=0$, every term with $p_X(x)>0$ must satisfy $\log(1/p_X(x))=0$, hence $p_X(x)=1$ on the support.
\end{proof}

\begin{theorem}[Maximum entropy on a finite alphabet]
If $|\cX|=M$, then $H(X)\le \log M$, with equality if and only if $X$ is uniform on $\cX$.
\end{theorem}
\begin{proof}
Let $U$ be the uniform pmf on $\cX$, so $U(x)=1/M$.
Using the KL divergence (defined in Section~\ref{sec:entropy-kl}) and its non-negativity,
\[
0\le \KL{p_X}{U}=\sum_x p_X(x)\log\frac{p_X(x)}{1/M} = -H(X)+\log M.
\]
Rearrange.
\end{proof}

\subsection{KL Divergence and Cross Entropy}
\label{sec:entropy-kl}

\begin{definition}[Kullback--Leibler (KL) divergence]
Let $P$ and $Q$ be pmfs on the same finite alphabet $\cX$.
The KL divergence from $P$ to $Q$ is
\begin{equation}
    \KL{P}{Q}
    \coloneqq \sum_{x:P(x)>0} P(x)\log\frac{P(x)}{Q(x)},
\end{equation}
with the convention that if $P(x)>0$ and $Q(x)=0$, then $\KL{P}{Q}=+\infty$.
\end{definition}

\paragraph{How to read $\KL{P}{Q}$.}
$\KL{P}{Q}$ is not symmetric and is not a metric.
Operationally, it measures the \emph{extra average code length} incurred when coding samples from $P$ using a code optimized for $Q$.

\begin{theorem}[Gibbs' inequality / KL non-negativity]
\label{thm:gibbs}
For pmfs $P,Q$ on the same finite alphabet,
\(\KL{P}{Q}\ge 0\), with equality if and only if $P=Q$ (with matching support).
\end{theorem}
\begin{proof}
Write
\[
-\KL{P}{Q} = \sum_x P(x)\log\frac{Q(x)}{P(x)} = \Esub{X\sim P}{\log\frac{Q(X)}{P(X)}}.
\]
Since $\log$ is concave, Jensen's inequality gives
\[
\Esub{X\sim P}{\log\frac{Q(X)}{P(X)}}
\le \log \Esub{X\sim P}{\frac{Q(X)}{P(X)}}
= \log \sum_x P(x)\frac{Q(x)}{P(x)}
= \log \sum_x Q(x)
= \log 1 =0.
\]
Thus $-\KL{P}{Q}\le 0$.
\end{proof}

\subsection{Cross Entropy Decomposition}
\label{sec:entropy-cross-entropy}

Let $X\sim P$ be the true distribution on $\cX$, and let $Q$ be a model distribution.
The expected log loss when using $Q$ is
\begin{equation}
    H(P,Q) \coloneqq \Esub{X\sim P}{\log\frac{1}{Q(X)}}
    = \sum_{x\in\cX} P(x)\log\frac{1}{Q(x)}.
\end{equation}
This quantity is often called the \emph{cross entropy} of $P$ relative to $Q$.

\begin{theorem}[Cross entropy decomposition]
\label{thm:cross-entropy-decomp}
For any pmfs $P,Q$ on the same alphabet,
\begin{equation}
    H(P,Q) = H(P) + \KL{P}{Q}.
\end{equation}
In particular, $H(P,Q)\ge H(P)$ with equality iff $P=Q$.
\end{theorem}
\begin{proof}
Expand and regroup:
\[
H(P,Q)-H(P)
= \sum_x P(x)\log\frac{1}{Q(x)} - \sum_x P(x)\log\frac{1}{P(x)}
= \sum_x P(x)\log\frac{P(x)}{Q(x)}
= \KL{P}{Q}.
\]
Apply Theorem~\ref{thm:gibbs}.
\end{proof}

\paragraph{Connection to classification.}
Section~\ref{sec:intro-logloss} showed that the standard log loss is a KL divergence.
Theorem~\ref{thm:cross-entropy-decomp} explains why minimizing cross entropy is the same as minimizing KL divergence up to an additive constant.

\subsection{Joint and Conditional Entropy}
\label{sec:entropy-conditional}

\begin{definition}[Joint entropy]
For jointly distributed discrete random variables $(X,Y)$ with joint pmf $p_{X,Y}$,
\begin{equation}
    H(X,Y)\coloneqq -\sum_{x,y} p_{X,Y}(x,y)\log p_{X,Y}(x,y).
\end{equation}
\end{definition}

\begin{definition}[Conditional entropy]
For $(X,Y)$ with conditional pmf $p_{X\mid Y}$, the conditional entropy of $X$ given $Y$ is
\begin{equation}
    H(X\mid Y)
    \coloneqq \E\!\left[\log\frac{1}{p_{X\mid Y}(X\mid Y)}\right]
    = \sum_{y} p_Y(y)\,H(X\mid Y=y).
\end{equation}
\end{definition}

\begin{theorem}[Chain rule for entropy]
\label{thm:entropy-chain-rule}
\begin{equation}
    H(X,Y) = H(Y) + H(X\mid Y).
\end{equation}
More generally, for $X^n=(X_1,\dots,X_n)$,
\begin{equation}
    H(X^n)=\sum_{i=1}^n H\bigl(X_i\mid X^{i-1}\bigr),\qquad X^{i-1}=(X_1,\dots,X_{i-1}).
\end{equation}
\end{theorem}
\begin{proof}
Use $p_{X,Y}(x,y)=p_Y(y)p_{X\mid Y}(x\mid y)$ and expand:
\[
\begin{aligned}
H(X,Y)
&=-\sum_{x,y}p_{X,Y}(x,y)\log\bigl(p_Y(y)p_{X\mid Y}(x\mid y)\bigr)\\
&=-\sum_{x,y}p_{X,Y}(x,y)\log p_Y(y)\; -\sum_{x,y}p_{X,Y}(x,y)\log p_{X\mid Y}(x\mid y)\\
&=H(Y)+H(X\mid Y).
\end{aligned}
\]
The multi-variable case follows by iterating the two-variable identity.
\end{proof}

\begin{theorem}[Conditioning reduces entropy]
\label{thm:conditioning-reduces-entropy}
For discrete $X,Y$,
\begin{equation}
    H(X\mid Y)\le H(X),
\end{equation}
with equality if and only if $X$ and $Y$ are independent.
\end{theorem}
\begin{proof}
Using Theorem~\ref{thm:entropy-chain-rule},
\(
H(X)-H(X\mid Y) = I(X;Y)
\)
where $I(X;Y)$ is mutual information (defined next) and is always nonnegative.
Independence is exactly the equality condition.
\end{proof}

\subsection{Mutual Information}
\label{sec:entropy-mi}

\begin{definition}[Mutual information]
\label{def:mi}
For jointly distributed discrete $(X,Y)$,
\begin{equation}
    I(X;Y)
    \coloneqq \KL{p_{X,Y}}{p_Xp_Y}
    = \sum_{x,y} p_{X,Y}(x,y)\log\frac{p_{X,Y}(x,y)}{p_X(x)p_Y(y)}.
\end{equation}
\end{definition}

\paragraph{Equivalent forms.}
Using the chain rule,
\begin{equation}
    I(X;Y)=H(X)-H(X\mid Y)=H(Y)-H(Y\mid X).
\end{equation}

\begin{theorem}[Basic properties]
\label{thm:mi-basic}
Mutual information satisfies:
\begin{enumerate}[leftmargin=2em]
    \item \textbf{Non-negativity:} $I(X;Y)\ge 0$, with equality iff $X$ and $Y$ are independent.
    \item \textbf{Symmetry:} $I(X;Y)=I(Y;X)$.
    \item \textbf{Upper bounds:} $I(X;Y)\le H(X)$ and $I(X;Y)\le H(Y)$.
\end{enumerate}
\end{theorem}
\begin{proof}
(1) is Theorem~\ref{thm:gibbs} applied to $\KL{p_{X,Y}}{p_Xp_Y}$.
(2) follows by swapping $X$ and $Y$ in the definition.
For (3), use $I(X;Y)=H(X)-H(X\mid Y)$ and $H(X\mid Y)\ge 0$.
\end{proof}

\begin{definition}[Conditional mutual information]
For $(X,Y,Z)$,
\begin{equation}
    I(X;Z\mid Y)
    \coloneqq H(X\mid Y)-H(X\mid Y,Z).
\end{equation}
Equivalently,
\begin{equation}
    I(X;Z\mid Y)=\Esub{Y}{\KL{p_{X,Z\mid Y}}{p_{X\mid Y}p_{Z\mid Y}}}\ge 0.
\end{equation}
\end{definition}

\begin{theorem}[Data processing inequality]
\label{thm:dpi}
If $X\to Y\to Z$ form a Markov chain (i.e., $X$ and $Z$ are conditionally independent given $Y$), then
\begin{equation}
    I(X;Z)\le I(X;Y).
\end{equation}
\end{theorem}
\begin{proof}
A standard identity is
\(I(X;Y,Z)=I(X;Y)+I(X;Z\mid Y)\).
Under the Markov condition, $I(X;Z\mid Y)=0$, so $I(X;Y,Z)=I(X;Y)$.
Finally, $I(X;Z)\le I(X;Y,Z)$ because observing $(Y,Z)$ cannot reveal less about $X$ than observing $Z$ alone.
\end{proof}

\begin{remark}[Why DPI matters]
DPI formalizes the idea that deterministic or stochastic post-processing cannot create information about the original source.
In deep learning, it is often invoked to reason about learned representations $X\to T\to \widehat Y$.
\end{remark}


\begin{remark}
Note that there are no ordering between $I(X;Z)$ and $I(X;Z|Y)$.
More interestingly, $I(X;Z)=0$ does not imply $I(X;Z|Y)=0$.
Consider the following example:
\[
X,Z\sim\text{Bernoulli}(1/2)\quad\text{and}\quad Y=X\oplus Z.
\]
Then:
\[
I(X;Z)=0,\quad\text{but}\quad I(X;Z\,|\,Y)=1,
\]
showing explicitly how conditioning can increase mutual information.
\end{remark}
